\chapter{Complete Code Reference}
\label{app:code}

This appendix provides a complete reference of all Python modules, functions, and classes available in the codebase.

\section{Module Overview}

\begin{table}[H]
\centering
\begin{tabular}{ll}
\toprule
\textbf{Module} & \textbf{Description} \\
\midrule
\texttt{modular.py} & Modular arithmetic, CRT, primitive roots, discrete log \\
\texttt{quadratic.py} & Legendre/Jacobi/Kronecker symbols, quadratic reciprocity \\
\texttt{gauss\_sums.py} & Gauss sums, Dirichlet characters, L-functions \\
\texttt{disquisitiones.py} & Binary quadratic forms, class groups, cyclotomic polynomials \\
\texttt{gaussian\_integral.py} & Error function, normal distribution, error propagation \\
\texttt{theta.py} & Jacobi theta functions, Dedekind eta, Eisenstein series \\
\texttt{agm.py} & Arithmetic-geometric mean, elliptic integrals, $\pi$ algorithms \\
\texttt{gaussian\_elim.py} & Gaussian elimination, LU/Cholesky/QR decomposition \\
\texttt{least\_squares.py} & Linear/nonlinear least squares, Gauss-Newton, Levenberg-Marquardt \\
\texttt{gauss\_quadrature.py} & Gauss-Legendre/Hermite/Laguerre/Chebyshev quadrature \\
\texttt{normal.py} & Normal distribution, MLE, confidence intervals \\
\texttt{gp.py} & Gaussian processes, Kriging, kernels \\
\texttt{heptadecagon.py} & 17-gon construction, Gauss periods, constructible polygons \\
\texttt{theorema.py} & Differential geometry, Theorema Egregium, Gauss-Bonnet \\
\texttt{orbital.py} & Orbital mechanics, Gauss's method, Kepler equation \\
\texttt{lemniscate.py} & Lemniscate functions, Gauss constant relations \\
\texttt{gauss\_composition.py} & Binary quadratic form composition, class groups \\
\bottomrule
\end{tabular}
\end{table}

\section{modular.py}

\begin{lstlisting}[style=python]
mod_inv(a, n)                              # Modular inverse
chinese_remainder(residues, moduli)        # Gauss's CRT
primitive_root(p)                          # Find primitive root mod p
all_primitive_roots(p)                     # All primitive roots
discrete_log_gauss(g, a, p)                # Baby-step giant-step
euler_totient(n)                           # Euler's phi function
multiplicative_order(a, n)                 # Order of a mod n
legendre_symbol(a, p)                      # Legendre symbol (a/p)
tonelli_shanks(n, p)                       # Square root mod p
\end{lstlisting}

\section{quadratic.py}

\begin{lstlisting}[style=python]
legendre_symbol(a, p)                      # Euler's criterion
legendre_symbol_gauss(a, p)                # Gauss's Lemma
jacobi_symbol(a, n)                        # Jacobi symbol
kronecker_symbol(a, n)                     # Kronecker symbol
quadratic_reciprocity(p, q)                # Verify (p/q)(q/p)
quadratic_residues(p)                      # List quadratic residues
quadratic_nonresidues(p)                   # List non-residues
tonelli_shanks(n, p)                       # Square root mod p
cipolla(a, p)                              # Cipolla's algorithm
solve_quadratic_congruence(a, b, c, p)     # ax^2+bx+c == 0 (mod p)
gauss_sum(p)                               # Quadratic Gauss sum
\end{lstlisting}

\section{gauss\_sums.py}

\begin{lstlisting}[style=python]
quadratic_gauss_sum(p)                     # g(p) = sum e^{2*pi*i*a^2/p}
general_gauss_sum(chi, n)                  # tau(chi) = sum chi(a) * e^{2*pi*i*a/n}
dirichlet_characters(n)                    # All characters mod n
primitive_character(n)                     # Find primitive character
dirichlet_l_function(chi, s)               # L(s, chi)
polya_vinogradov(chi, N)                   # Character sum bound
verify_gauss_sum(p)                        # Check g(p) = sqrt(p) or i*sqrt(p)
\end{lstlisting}

\section{disquisitiones.py}

\begin{lstlisting}[style=python]
solve_linear_congruence(a, b, n)           # ax == b (mod n)
is_primitive(a, b, c)                      # Check if form is primitive
discriminant(a, b, c)                       # Delta = b^2 - 4ac
reduced_forms(D)                           # All reduced forms of disc D
class_number(D)                            # h(D)
form_multiply(f1, f2)                      # Gauss composition
class_group(D)                             # Class group as list of forms
class_group_structure(D)                   # Group invariants
cyclotomic_poly(n)                         # Phi_n(x) coefficients
primitive_root(p)                          # Primitive root mod p
heptadecagon_vertices()                    # 17-gon coordinates
fermat_primes()                            # Known Fermat primes
is_constructible(n)                        # Check n-gon constructible
represents(a, b, c, n)                     # Does form represent n?
\end{lstlisting}

\section{gaussian\_integral.py}

\begin{lstlisting}[style=python]
gaussian_integral(a, b, c, x1, x2)         # integral exp(-ax^2+bx+c) dx
gaussian_pdf(x, mu, sigma)                 # Normal PDF
gaussian_cdf(x, mu, sigma)                 # Normal CDF
gaussian_ppf(p, mu, sigma)                 # Normal quantile
erf(x)                                     # Error function
erfc(x)                                    # Complementary erf
erfi(x)                                    # Imaginary error function
dawson_integral(x)                         # Dawson's integral
faddeeva(z)                                # Faddeeva function w(z)
normal_pdf(x, mu, sigma)                   # Normal PDF
normal_cdf(x, mu, sigma)                   # Normal CDF
normal_sf(x, mu, sigma)                    # Survival function
normal_ppf(p, mu, sigma)                   # Normal quantile
normal_rvs(size, mu, sigma)                # Random variates
confidence_interval_normal(mean, std, n)   # CI for mean
error_propagation(f, vars)                 # Delta method
chi2_pdf(x, df), t_pdf(x, df), f_pdf(x, df1, df2)
\end{lstlisting}

\section{theta.py}

\begin{lstlisting}[style=python]
theta1(z, tau), theta2(z, tau)             # Jacobi theta functions
theta3(z, tau), theta4(z, tau)
theta_null(tau)                            # theta(0|tau) tuple
dedekind_eta(tau)                          # Dedekind eta function
theta_quadratic_form(A, z)                 # Theta of quadratic form
sum_of_squares(k, N)                       # r_k(n) for n<=N
jacobi_triple_product(q, z)                # Identity for theta_1
eisenstein_series(k, tau)                  # G_k(tau)
j_invariant(tau)                           # j(tau)
modular_transform_sl2z(tau, a, b, c, d)    # SL(2,Z) action
\end{lstlisting}

\section{agm.py}

\begin{lstlisting}[style=python]
agm(a, b)                                  # Arithmetic-geometric mean
agm_sequence(a, b)                         # Full iteration sequence
gauss_constant()                           # 1/M(1, sqrt(2))
elliptic_K(k)                              # Complete elliptic integral K
elliptic_E(k)                              # Complete elliptic integral E
pi_brent_salamin(n)                        # pi via Brent-Salamin
pi_borwein_quartic(n)                      # pi via Borwein quartic
lemniscate_constant()                      # Arc length of lemniscate
elliptic_perimeter(a, b)                   # Ellipse perimeter
theta_null_via_agm(tau)                    # theta_3(0|tau) via AGM
hypergeometric_2F1(a, b, c, z)             # Gauss hypergeometric
agm_hypergeometric(a, b)                   # M(a,b) via 2F1
\end{lstlisting}

\section{gaussian\_elim.py}

\begin{lstlisting}[style=python]
gaussian_elimination(A, b)                 # Solve Ax=b
lu_decomposition(A)                        # PA=LU
lu_solve(P, L, U, b)                       # Solve PAx=b
cholesky(A)                                # Cholesky LL^T
cholesky_solve(L, b)                       # Solve LL^Tx=b
qr_decomposition(A)                        # QR via Householder
qr_solve(A, b)                             # Least squares via QR
matrix_inverse(A)                          # Inverse via LU
condition_number(A, norm)                  # kappa(A)
least_squares_normal(A, b)                 # Normal equations
least_squares_qr(A, b)                     # QR least squares
ridge_regression(A, b, lam)                # Tikhonov regularization
bareiss(A)                                 # Fraction-free elimination
\end{lstlisting}

\section{least\_squares.py}

\begin{lstlisting}[style=python]
least_squares_normal(A, b)                 # Normal equations
least_squares_qr(A, b)                     # QR least squares
weighted_least_squares(A, b, W)            # Weighted LS
ridge_regression(A, b, lam)                # Ridge regression
gauss_newton(f, jac, x0)                   # Gauss-Newton
levenberg_marquardt(f, jac, x0)            # Levenberg-Marquardt
fit_exponential(x, y)                      # a*exp(bx)+c fit
polyfit(x, y, deg)                         # Polynomial fit
polyval(coeffs, x)                         # Evaluate polynomial
ceres_orbit(observations)                  # Gauss's Ceres method
\end{lstlisting}

\section{gauss\_quadrature.py}

\begin{lstlisting}[style=python]
gauss_legendre(n)                          # Nodes/weights for [-1,1]
gauss_hermite(n)                           # (-inf, inf) weight e^{-x^2}
gauss_laguerre(n, alpha)                   # [0, inf) weight x^alpha * e^{-x}
gauss_chebyshev(n, kind)                   # [-1,1] weight (1-x^2)^{+/-1/2}
golub_welsch(alpha, beta)                  # General Golub-Welsch
integrate(f, a, b, n)                      # Quadrature on [a,b]
integrate_infinite(f, n)                   # Quadrature on (-inf, inf)
integrate_adaptive(f, a, b)                # Adaptive Gauss-Kronrod
gauss_kronrod(n)                           # Gauss-Kronrod nodes/weights
\end{lstlisting}

\section{normal.py}

\begin{lstlisting}[style=python]
normal_pdf(x, mu, sigma)                   # Normal PDF
normal_cdf(x, mu, sigma)                   # Normal CDF
normal_sf(x, mu, sigma)                    # Survival function
normal_ppf(p, mu, sigma)                   # Quantile function
normal_rvs(size, mu, sigma)                # Random variates
normal_mle(data)                           # MLE for mu, sigma
normal_mle_unbiased(data)                  # Unbiased estimators
confidence_interval_normal(mean, std, n)   # CI for mean
confidence_interval_t(mean, std, n)        # t-distribution CI
z_test(data, mu0, sigma)                   # Z-test
t_test(data, mu0)                          # t-test
error_propagation(f, vars)                 # Delta method
multivariate_normal_pdf(x, mu, Sigma)      # Multivariate PDF
multivariate_normal_logpdf(x, mu, Sigma)   # Multivariate log PDF
chi2_pdf(x, df), chi2_cdf(x, df)
t_pdf(x, df), f_pdf(x, df1, df2)
box_muller(u1, u2)                         # Box-Muller transform
marsaglia_bray()                           # Polar method
\end{lstlisting}

\section{gp.py}

\begin{lstlisting}[style=python]
class GaussianProcess:
    __init__(kernel, noise=1e-6)
    fit(X, y, **kernel_params)
    predict(X_star, return_std=True)
    log_marginal_likelihood(**kernel_params)
    optimize_hyperparameters(X, y, initial_params, bounds)

rbf_kernel(x1, x2, length_scale, variance)
matern_kernel(x1, x2, length_scale, variance, nu)
periodic_kernel(x1, x2, length_scale, variance, period)
rational_quadratic_kernel(x1, x2, length_scale, variance, alpha)
linear_kernel(x1, x2, variance, c)
white_kernel(x1, x2, noise_level)
simple_kriging(x_obs, y_obs, x_pred, kernel)
ordinary_kriging(x_obs, y_obs, x_pred, kernel)
\end{lstlisting}

\section{heptadecagon.py}

\begin{lstlisting}[style=python]
cyclotomic_poly(n)                         # Phi_n(x) coefficients
primitive_root(p)                          # Primitive root mod p
constructible_ngons(limit)                 # List constructible n
heptadecagon_vertices()                    # 17-gon coordinates
cos_2pi_over_17()                          # Exact expression
cos_2pi_over_17_value()                    # Numerical value
gaussian_periods(p, length)                # Gauss periods
fermat_primes()                            # [3, 5, 17, 257, 65537]
is_constructible(n)                        # Check n-gon constructible
\end{lstlisting}

\section{theorema.py}

\begin{lstlisting}[style=python]
first_fundamental_form(r, u, v)            # E, F, G
second_fundamental_form(r, u, v)           # L, M, N
gaussian_curvature(E, F, G, L, M, N)       # K = (LN-M^2)/(EG-F^2)
mean_curvature(E, F, G, L, M, N)           # H = (EN-2FM+GL)/(2(EG-F^2))
principal_curvatures(E, F, G, L, M, N)     # k_1, k_2
christoffel_symbols(E, F, G, ...)          # Gamma^k_ij
christoffel_from_surface(r, u, v)          # From parametrization
theorema_egregium_check(r, u, v)           # Verify K intrinsic
geodesic_equations(E, F, G, ...)           # Geodesic ODEs
geodesic_ode(state, s, r)                  # Geodesic ODE system
gauss_bonnet_discrete(vertices, faces)     # Discrete Gauss-Bonnet

# Example surfaces
sphere(u, v, R)
cylinder(u, v, R)
torus(u, v, R, r)
pseudosphere(u, v)
\end{lstlisting}

\section{orbital.py}

\begin{lstlisting}[style=python]
solve_kepler(M, e)                         # Kepler's equation E - e sin E = M
mean_anomaly_from_true(f, e)               # f -> M
true_anomaly_from_mean(M, e)               # M -> f
state_to_elements(r, v, mu)                # State -> orbital elements
elements_to_state(elems, mu)               # Elements -> state vector
rotation_matrix(Omega, i, omega)           # Perifocal -> inertial
fg_functions(r0, v0, dt, mu)               # f, g, fdot, gdot
universal_anomaly(r0, rv0, alpha, dt, mu)  # Universal anomaly chi
stumpff_C(z), stumpff_S(z)                 # Stumpff functions
propagate_universal(r0, v0, dt, mu)        # Universal propagation
gauss_method(obs1, obs2, obs3, mu)         # Gauss's 3-obs method
differential_correction(obs, x0, mu)       # Batch least squares
ceres_example()                            # Ceres orbit demo
\end{lstlisting}

\section{lemniscate.py}

\begin{lstlisting}[style=python]
LEMNISCATE_CONSTANT                        # omega_bar ~ 2.622057554
GAUSS_CONSTANT                             # G ~ 0.834626842
lemniscate_integral(u)                     # integral_0^u dt/sqrt(1-t^4)
lemniscate_integral_inverse(omega)         # sl(omega) = u
sl(omega)                                  # Lemniscate sine
cl(omega)                                  # Lemniscate cosine
pythagorean_check(omega)                   # sl^2+cl^2+sl^2*cl^2
sl_add(u, v)                               # sl(u+v) formula
cl_add(u, v)                               # cl(u+v) formula
lemniscate_division(n)                     # n-division points
constructible_lemniscate_divisions(limit)  # Constructible n
lemniscate_hypergeometric(x)               # omega = x * 2F1(1/4, 1/2; 5/4; x^4)
\end{lstlisting}

\section{gauss\_composition.py}

\begin{lstlisting}[style=python]
class BinaryQuadraticForm:
    __init__(a, b, c)                      # ax^2 + bxy + cy^2
    is_primitive()                         # Check gcd=1
    is_reduced()                           # Check Gauss reduction
    value(x, y)                            # Evaluate form
    __repr__()                             # String representation

discriminant(a, b, c)                       # Delta = b^2 - 4ac
reduce_form(f)                             # Gauss reduction
all_reduced_forms(D)                       # All reduced forms
class_number(D)                            # h(D)
gauss_composition(f, g)                    # Gauss composition
compose_forms(f, g)                        # Alias

class ClassGroup:
    __init__(D)                            # Discriminant D
    compose(f, g)                          # Group operation
    inverse(f)                             # Inverse element
    order(f)                               # Order of element
    is_cyclic()                            # Check if cyclic
    invariants()                           # Group structure

gauss_class_numbers(limit)                 # Class numbers for D<0
class_number_one_discriminants()           # D with h(D)=1
verify_gauss_class_numbers()               # Check against Gauss's table
dirichlet_composition(f, g)                # Dirichlet's algorithm
mobius(n)                                  # Mobius function
\end{lstlisting}

\section{Quick Start Examples}

\subsection{Chapter 1: Modular Arithmetic}
\begin{lstlisting}[style=python]
from python.modular import chinese_remainder, primitive_root, discrete_log_gauss

# CRT: x ≡ 2 (mod 3), x ≡ 3 (mod 5), x ≡ 2 (mod 7)
x = chinese_remainder([2, 3, 2], [3, 5, 7])  # x = 23

# Primitive root mod 17
g = primitive_root(17)  # 3

# Discrete log: 3^x ≡ 13 (mod 17)
x = discrete_log_gauss(3, 13, 17)  # x = 4
\end{lstlisting}

\subsection{Chapter 2: Quadratic Reciprocity}
\begin{lstlisting}[style=python]
from python.quadratic import legendre_symbol, jacobi_symbol, tonelli_shanks

# Legendre symbol
print(legendre_symbol(5, 13))  # 1 (5 is QR mod 13)

# Quadratic reciprocity
print(legendre_symbol(5, 13) * legendre_symbol(13, 5))  # 1

# Square root mod p
r = tonelli_shanks(10, 13)  # 6 or 7 (since 6² = 36 ≡ 10 mod 13)
\end{lstlisting}

\subsection{Chapter 7: AGM and \texorpdfstring{$\pi$}{Pi}}
\begin{lstlisting}[style=python]
from python.agm import agm, pi_brent_salamin, elliptic_K

# AGM
print(agm(1, 0.5))  # 0.7283958933...

# pi via Brent-Salamin
for n in range(1, 6):
    print(pi_brent_salamin(n))  # Converges quadratically

# Elliptic integral
K = elliptic_K(1/math.sqrt(2))  # Gamma(1/4)^2 / (4*sqrt(pi))
\end{lstlisting}

\subsection{Chapter 8: Gaussian Elimination}
\begin{lstlisting}[style=python]
from python.gaussian_elim import gaussian_elimination, lu_decomposition, qr_solve
import numpy as np

A = np.array([[2, 1, -1], [-3, -1, 2], [-2, 1, 2]], dtype=float)
b = np.array([8, -11, -3], dtype=float)

x = gaussian_elimination(A, b)  # [2, 3, -1]
P, L, U = lu_decomposition(A)
x_qr = qr_solve(A, b)  # Least squares via QR
\end{lstlisting}

\subsection{Chapter 13: Heptadecagon}
\begin{lstlisting}[style=python]
from python.heptadecagon import heptadecagon_vertices, constructible_ngons

verts = heptadecagon_vertices()
print(verts[:3])  # First 3 vertices

print(constructible_ngons(50))
# [3, 4, 5, 6, 8, 10, 12, 15, 16, 17, 20, 24, 30, 32, 34, 40, 48]
\end{lstlisting}

\subsection{Chapter 14: Theorema Egregium}
\begin{lstlisting}[style=python]
from python.theorema import sphere, gaussian_curvature, first_fundamental_form

r = lambda u, v: (math.sin(u)*math.cos(v), math.sin(u)*math.sin(v), math.cos(u))
E, F, G = first_fundamental_form(r, 1.0, 0.5)
L, M, N = second_fundamental_form(r, 1.0, 0.5)
K = gaussian_curvature(E, F, G, L, M, N)  # K = 1.0 for unit sphere
\end{lstlisting}

\section{Version Information}

All modules are part of the \textit{Exploring Gauss's Mathematics using Python} codebase, version 1.0.

For updates and contributions, visit the project repository.